\documentclass[10pt,a4paper,full]{physicsquiz}

\quiztitle{PHY 104 Calculation-Only Question Bank}
\quizcourse{100-Question Numerical Examination Practice}
\quizauthor{Ot}
\quizinstitution{}
\quizversion{1.0}
\quizinstructions{Choose the best numerical answer. Show working before selecting an option. All questions carry one mark. A calculator may be used.}
\quizconstants{
  \(g=\SI{9.80}{m.s^{-2}}\),
  \(v_{\text{sound, air}}=\SI{343}{m.s^{-1}}\),
  \(c=\SI{3.00e8}{m.s^{-1}}\), and
  \(I_{\mathrm{ref}}=\SI{1.00e-12}{W.m^{-2}}\)
}

\newcommand{\fivechoices}[5]{%
  \item #1
  \item #2
  \item #3
  \item #4
  \item #5
}

\newcommand{\bankq}[9]{%
  \begin{quizquestion}[
    id=#1,
    topic={#2},
    difficulty=#3,
    marks=1,
    correct=#4,
    tags={#5},
    outcome={#6}
  ]
    #7
    \begin{choiceoptions}
      #8
    \end{choiceoptions}
  \end{quizquestion}
  \begin{quizsolution}
    #9
  \end{quizsolution}
}

\begin{document}
\hypersetup{pageanchor=false}
\makequiztitle
\hypersetup{pageanchor=true}
\pagenumbering{arabic}

\begin{quizreferencecontent}
\section*{Calculation Coverage}
\begin{multicols}{2}
\begin{itemize}[leftmargin=*]
  \item SHM kinematics and energy
  \item Damping, resonance, transients, and quality factor
  \item Coupled oscillators and normal modes
  \item Travelling waves, power, and intensity
  \item Interference, beats, and optical path calculations
  \item Diffraction and polarization
  \item Phase velocity, group velocity, and dispersion
  \item Wave equations and bounded modes in 1D, 2D, and 3D
  \item Echoes, pipes, sound levels, and Doppler shift
  \item Propagation, impedance, transmission, and refraction
\end{itemize}
\end{multicols}
\end{quizreferencecontent}

\begin{quizbank}

% 1--10: SHM and energy (3 foundation, 5 applied, 2 challenge)
\bankq{calc-shm-001}{SHM Kinematics and Energy}{foundation}{A}{shm,period}{Calculate a mass-spring period.}{%
A \(\SI{0.50}{kg}\) mass is attached to a spring of constant \(\SI{200}{N.m^{-1}}\). What is its period?}{%
\fivechoices{\(\SI{0.314}{s}\)}{\(\SI{0.444}{s}\)}{\(\SI{0.628}{s}\)}{\(\SI{1.57}{s}\)}{\(\SI{2.00}{s}\)}}{%
\(T=2\pi\sqrt{m/k}=2\pi\sqrt{0.50/200}=\SI{0.314}{s}\).}

\bankq{calc-shm-002}{SHM Kinematics and Energy}{foundation}{B}{shm,speed}{Calculate maximum SHM speed.}{%
An oscillator has amplitude \(\SI{0.080}{m}\) and frequency \(\SI{2.0}{Hz}\). Find its maximum speed.}{%
\fivechoices{\(\SI{0.503}{m.s^{-1}}\)}{\(\SI{1.01}{m.s^{-1}}\)}{\(\SI{1.26}{m.s^{-1}}\)}{\(\SI{2.01}{m.s^{-1}}\)}{\(\SI{4.02}{m.s^{-1}}\)}}{%
\(v_{\max}=\omega A=2\pi fA=2\pi(2.0)(0.080)=\SI{1.01}{m.s^{-1}}\).}

\bankq{calc-shm-003}{SHM Kinematics and Energy}{foundation}{C}{shm,energy}{Calculate total SHM energy.}{%
A spring of constant \(\SI{250}{N.m^{-1}}\) oscillates with amplitude \(\SI{0.120}{m}\). Find the total energy.}{%
\fivechoices{\(\SI{0.90}{J}\)}{\(\SI{1.20}{J}\)}{\(\SI{1.80}{J}\)}{\(\SI{3.60}{J}\)}{\(\SI{7.20}{J}\)}}{%
\(E=\tfrac12kA^2=\tfrac12(250)(0.120)^2=\SI{1.80}{J}\).}

\bankq{calc-shm-004}{SHM Kinematics and Energy}{applied}{D}{shm,kinematics}{Calculate speed at a specified displacement.}{%
An oscillator has \(A=\SI{0.100}{m}\), \(\omega=\SI{8.0}{rad.s^{-1}}\), and instantaneous displacement \(x=\SI{0.060}{m}\). Find its speed.}{%
\fivechoices{\(\SI{0.32}{m.s^{-1}}\)}{\(\SI{0.48}{m.s^{-1}}\)}{\(\SI{0.60}{m.s^{-1}}\)}{\(\SI{0.64}{m.s^{-1}}\)}{\(\SI{0.80}{m.s^{-1}}\)}}{%
\(v=\omega\sqrt{A^2-x^2}=8\sqrt{0.100^2-0.060^2}=\SI{0.64}{m.s^{-1}}\).}

\bankq{calc-shm-005}{SHM Kinematics and Energy}{applied}{E}{shm,acceleration}{Calculate SHM acceleration from period and displacement.}{%
An oscillator of period \(\SI{0.500}{s}\) is at \(x=\SI{0.040}{m}\). What is the magnitude of its acceleration?}{%
\fivechoices{\(\SI{1.58}{m.s^{-2}}\)}{\(\SI{3.16}{m.s^{-2}}\)}{\(\SI{4.00}{m.s^{-2}}\)}{\(\SI{5.03}{m.s^{-2}}\)}{\(\SI{6.32}{m.s^{-2}}\)}}{%
\(\omega=2\pi/T=4\pi\), so \(|a|=\omega^2|x|=(4\pi)^2(0.040)=\SI{6.32}{m.s^{-2}}\).}

\bankq{calc-shm-006}{SHM Kinematics and Energy}{applied}{A}{pendulum,period}{Calculate pendulum length from period.}{%
A simple pendulum has period \(\SI{2.00}{s}\). Find its length.}{%
\fivechoices{\(\SI{0.993}{m}\)}{\(\SI{1.24}{m}\)}{\(\SI{1.56}{m}\)}{\(\SI{2.48}{m}\)}{\(\SI{9.80}{m}\)}}{%
\(L=g(T/2\pi)^2=9.80(2.00/2\pi)^2=\SI{0.993}{m}\).}

\bankq{calc-shm-007}{SHM Kinematics and Energy}{applied}{B}{shm,energy}{Calculate amplitude from energy.}{%
An oscillator has total energy \(\SI{0.720}{J}\) and spring constant \(\SI{100}{N.m^{-1}}\). Find its amplitude.}{%
\fivechoices{\(\SI{0.085}{m}\)}{\(\SI{0.120}{m}\)}{\(\SI{0.144}{m}\)}{\(\SI{0.240}{m}\)}{\(\SI{0.720}{m}\)}}{%
\(A=\sqrt{2E/k}=\sqrt{2(0.720)/100}=\SI{0.120}{m}\).}

\bankq{calc-shm-008}{SHM Kinematics and Energy}{applied}{C}{shm,phase}{Calculate phase from displacement and velocity sign.}{%
For \(x=A\cos\phi\), an oscillator is at \(x=A/2\) and moving toward decreasing \(x\). What is the smallest positive phase \(\phi\)?}{%
\fivechoices{\(\pi/6\)}{\(\pi/4\)}{\(\pi/3\)}{\(2\pi/3\)}{\(5\pi/3\)}}{%
\(\cos\phi=1/2\). Since \(v=-A\omega\sin\phi<0\), \(\sin\phi>0\), hence \(\phi=\pi/3\).}

\bankq{calc-shm-009}{SHM Kinematics and Energy}{challenge}{D}{shm,time}{Calculate first-passage time in SHM.}{%
An oscillator obeys \(x=A\cos(2.5t)\) in SI units. How long after \(t=0\) does it first pass through equilibrium?}{%
\fivechoices{\(\SI{0.157}{s}\)}{\(\SI{0.314}{s}\)}{\(\SI{0.500}{s}\)}{\(\SI{0.628}{s}\)}{\(\SI{1.257}{s}\)}}{%
At the first equilibrium crossing, \(2.5t=\pi/2\). Thus \(t=\pi/5=\SI{0.628}{s}\).}

\bankq{calc-shm-010}{SHM Kinematics and Energy}{challenge}{E}{spring,scaling}{Calculate period after changing spring length.}{%
A uniform spring gives a mass a period of \(\SI{0.800}{s}\). The spring is cut in half and the same mass is attached to one half. Find the new period.}{%
\fivechoices{\(\SI{0.283}{s}\)}{\(\SI{0.400}{s}\)}{\(\SI{0.466}{s}\)}{\(\SI{0.520}{s}\)}{\(\SI{0.566}{s}\)}}{%
Halving the spring length doubles \(k\). Therefore \(T'=T/\sqrt2=0.800/\sqrt2=\SI{0.566}{s}\).}

% 11--20: Damping, resonance, and transients (3 foundation, 5 applied, 2 challenge)
\bankq{calc-damp-011}{Damping, Resonance, and Transients}{foundation}{A}{damping,decay}{Calculate the damping constant.}{%
For \(m\ddot x+b\dot x+kx=0\), let \(m=\SI{0.30}{kg}\) and \(b=\SI{1.20}{kg.s^{-1}}\). Find \(\gamma=b/(2m)\).}{%
\fivechoices{\(\SI{2.00}{s^{-1}}\)}{\(\SI{1.20}{s^{-1}}\)}{\(\SI{0.600}{s^{-1}}\)}{\(\SI{0.250}{s^{-1}}\)}{\(\SI{4.00}{s^{-1}}\)}}{%
\(\gamma=b/(2m)=1.20/[2(0.30)]=\SI{2.00}{s^{-1}}\).}

\bankq{calc-damp-012}{Damping, Resonance, and Transients}{foundation}{B}{damping,frequency}{Calculate the damped natural frequency.}{%
An oscillator has \(\omega_0=\SI{10.0}{rad.s^{-1}}\) and \(\gamma=\SI{3.0}{s^{-1}}\). Find its damped natural frequency.}{%
\fivechoices{\(\SI{7.00}{rad.s^{-1}}\)}{\(\SI{9.54}{rad.s^{-1}}\)}{\(\SI{10.4}{rad.s^{-1}}\)}{\(\SI{13.0}{rad.s^{-1}}\)}{\(\SI{91.0}{rad.s^{-1}}\)}}{%
\(\omega'=\sqrt{\omega_0^2-\gamma^2}=\sqrt{100-9}=\SI{9.54}{rad.s^{-1}}\).}

\bankq{calc-damp-013}{Damping, Resonance, and Transients}{foundation}{C}{damping,half-life}{Calculate amplitude half-life.}{%
The amplitude envelope is \(A(t)=A_0e^{-0.40t}\). Find the time for the amplitude to fall to \(A_0/2\).}{%
\fivechoices{\(\SI{0.693}{s}\)}{\(\SI{1.00}{s}\)}{\(\SI{1.73}{s}\)}{\(\SI{2.50}{s}\)}{\(\SI{4.33}{s}\)}}{%
\(1/2=e^{-0.40t}\), so \(t=\ln2/0.40=\SI{1.73}{s}\).}

\bankq{calc-damp-014}{Damping, Resonance, and Transients}{applied}{D}{damping,critical}{Calculate critical damping coefficient.}{%
A mass-spring system has \(m=\SI{0.50}{kg}\) and \(k=\SI{72}{N.m^{-1}}\). Find the critical damping coefficient.}{%
\fivechoices{\(\SI{3.0}{kg.s^{-1}}\)}{\(\SI{6.0}{kg.s^{-1}}\)}{\(\SI{8.5}{kg.s^{-1}}\)}{\(\SI{12.0}{kg.s^{-1}}\)}{\(\SI{36.0}{kg.s^{-1}}\)}}{%
\(b_c=2\sqrt{mk}=2\sqrt{(0.50)(72)}=\SI{12.0}{kg.s^{-1}}\).}

\bankq{calc-damp-015}{Damping, Resonance, and Transients}{applied}{E}{resonance,amplitude}{Calculate resonant amplitude.}{%
A lightly damped oscillator has force amplitude \(F_0=\SI{6.0}{N}\), damping coefficient \(b=\SI{0.30}{kg.s^{-1}}\), and \(\omega_0=\SI{20}{rad.s^{-1}}\). Estimate its resonant amplitude.}{%
\fivechoices{\(\SI{0.10}{m}\)}{\(\SI{0.25}{m}\)}{\(\SI{0.50}{m}\)}{\(\SI{0.75}{m}\)}{\(\SI{1.00}{m}\)}}{%
For light damping, \(X_{\max}\approx F_0/(b\omega_0)=6.0/[0.30(20)]=\SI{1.00}{m}\).}

\bankq{calc-damp-016}{Damping, Resonance, and Transients}{applied}{A}{quality-factor,damping}{Calculate quality factor.}{%
An oscillator has \(m=\SI{0.80}{kg}\), \(\omega_0=\SI{25}{rad.s^{-1}}\), and \(b=\SI{0.50}{kg.s^{-1}}\). Find \(Q\).}{%
\fivechoices{40}{20}{16}{10}{5}{}}{%
\(Q=m\omega_0/b=(0.80)(25)/0.50=40\).}

\bankq{calc-damp-017}{Damping, Resonance, and Transients}{applied}{B}{quality-factor,bandwidth}{Calculate resonance bandwidth.}{%
A resonance has central frequency \(\SI{120}{Hz}\) and \(Q=30\). Find its bandwidth.}{%
\fivechoices{\(\SI{2.0}{Hz}\)}{\(\SI{4.0}{Hz}\)}{\(\SI{6.0}{Hz}\)}{\(\SI{15}{Hz}\)}{\(\SI{30}{Hz}\)}}{%
\(\Delta f=f_0/Q=120/30=\SI{4.0}{Hz}\).}

\bankq{calc-damp-018}{Damping, Resonance, and Transients}{applied}{C}{resonance,power}{Calculate average resonant power.}{%
A sinusoidal force of amplitude \(\SI{4.0}{N}\) drives an oscillator with \(b=\SI{0.50}{kg.s^{-1}}\) at resonance. Find the average input power.}{%
\fivechoices{\(\SI{4}{W}\)}{\(\SI{8}{W}\)}{\(\SI{16}{W}\)}{\(\SI{24}{W}\)}{\(\SI{32}{W}\)}}{%
At resonance, \(\overline P=F_0^2/(2b)=4.0^2/[2(0.50)]=\SI{16}{W}\).}

\bankq{calc-damp-019}{Damping, Resonance, and Transients}{challenge}{D}{resonance,phase}{Calculate phase lag of a driven oscillator.}{%
For \(m=\SI{1.0}{kg}\), \(k=\SI{100}{N.m^{-1}}\), \(b=\SI{4.0}{kg.s^{-1}}\), and drive frequency \(\omega=\SI{8.0}{rad.s^{-1}}\), find the phase lag \(\delta\).}{%
\fivechoices{\(20.8^\circ\)}{\(28.1^\circ\)}{\(36.9^\circ\)}{\(41.6^\circ\)}{\(48.4^\circ\)}}{%
\(\tan\delta=b\omega/(k-m\omega^2)=32/(100-64)=0.8889\), so \(\delta=41.6^\circ\).}

\bankq{calc-damp-020}{Damping, Resonance, and Transients}{challenge}{E}{transient,beats}{Calculate transient beating time.}{%
A transient contains angular frequencies \(\SI{12.0}{rad.s^{-1}}\) and \(\SI{11.2}{rad.s^{-1}}\). Find the envelope period.}{%
\fivechoices{\(\SI{0.80}{s}\)}{\(\SI{1.25}{s}\)}{\(\SI{3.93}{s}\)}{\(\SI{6.28}{s}\)}{\(\SI{7.85}{s}\)}}{%
The envelope period is \(2\pi/|\Delta\omega|=2\pi/0.80=\SI{7.85}{s}\).}

% 21--30: Coupled SHM and normal modes (3 foundation, 4 applied, 3 challenge)
\bankq{calc-mode-021}{Coupled SHM and Normal Modes}{foundation}{A}{normal-modes,frequency}{Calculate the in-phase normal frequency.}{%
Two identical masses \(m=\SI{2.0}{kg}\) are attached to end springs of constant \(k=\SI{50}{N.m^{-1}}\) and coupled symmetrically. Find the in-phase angular frequency.}{%
\fivechoices{\(\SI{5.00}{rad.s^{-1}}\)}{\(\SI{7.07}{rad.s^{-1}}\)}{\(\SI{10.0}{rad.s^{-1}}\)}{\(\SI{25.0}{rad.s^{-1}}\)}{\(\SI{50.0}{rad.s^{-1}}\)}}{%
The coupling spring is unstretched in phase, so \(\omega_+=\sqrt{k/m}=\sqrt{50/2}=\SI{5.00}{rad.s^{-1}}\).}

\bankq{calc-mode-022}{Coupled SHM and Normal Modes}{foundation}{B}{normal-modes,frequency}{Calculate the out-of-phase normal frequency.}{%
For the system in Question 21, let the coupling spring constant be \(k_c=\SI{25}{N.m^{-1}}\). Find the out-of-phase angular frequency.}{%
\fivechoices{\(\SI{5.00}{rad.s^{-1}}\)}{\(\SI{7.07}{rad.s^{-1}}\)}{\(\SI{8.66}{rad.s^{-1}}\)}{\(\SI{10.0}{rad.s^{-1}}\)}{\(\SI{12.5}{rad.s^{-1}}\)}}{%
\(\omega_- =\sqrt{(k+2k_c)/m}=\sqrt{(50+50)/2}=\sqrt{50}=\SI{7.07}{rad.s^{-1}}\).}

\bankq{calc-mode-023}{Coupled SHM and Normal Modes}{foundation}{C}{coupled-oscillators,energy-transfer}{Calculate energy-transfer time from mode splitting.}{%
Two coupled modes have angular frequencies \(\SI{6.0}{rad.s^{-1}}\) and \(\SI{8.0}{rad.s^{-1}}\). Estimate the first complete energy-transfer time.}{%
\fivechoices{\(\SI{0.785}{s}\)}{\(\SI{1.00}{s}\)}{\(\SI{1.57}{s}\)}{\(\SI{3.14}{s}\)}{\(\SI{6.28}{s}\)}}{%
The first transfer occurs at \(t=\pi/(\omega_2-\omega_1)=\pi/2=\SI{1.57}{s}\).}

\bankq{calc-mode-024}{Coupled SHM and Normal Modes}{applied}{D}{eigenvalues,normal-modes}{Calculate a normal frequency from a stiffness matrix.}{%
For unit mass, the stiffness matrix is \(K=\begin{pmatrix}5&-1\\-1&5\end{pmatrix}\). Find the higher normal angular frequency.}{%
\fivechoices{\(\sqrt2\)}{2}{\(\sqrt5\)}{\(\sqrt6\)}{3}{}}{%
The eigenvalues of \(K\) are \(4\) and \(6\), so the higher frequency is \(\omega=\sqrt6\).}

\bankq{calc-mode-025}{Coupled SHM and Normal Modes}{applied}{E}{normal-modes,normalization}{Normalize a mode vector.}{%
The low-frequency mode is proportional to \([1,1]^T\). What is the first component of its unit-normalized vector?}{%
\fivechoices{0.250}{0.500}{0.577}{0.667}{0.707}{}}{%
The norm is \(\sqrt{1^2+1^2}=\sqrt2\), so the component is \(1/\sqrt2=0.707\).}

\bankq{calc-mode-026}{Coupled SHM and Normal Modes}{applied}{A}{normal-modes,decomposition}{Calculate modal amplitudes.}{%
Write \([3,1]^T=a[1,1]^T+b[1,-1]^T\). Find \((a,b)\).}{%
\fivechoices{\((2,1)\)}{\((1,2)\)}{\((3,1)\)}{\((2,-1)\)}{\((4,2)\)}}{%
The component equations are \(a+b=3\) and \(a-b=1\). Hence \(a=2\), \(b=1\).}

\bankq{calc-mode-027}{Coupled SHM and Normal Modes}{applied}{B}{coupled-oscillators,beats}{Calculate transfer time from two frequencies.}{%
Weakly coupled oscillators have normal angular frequencies \(\SI{4.0}{rad.s^{-1}}\) and \(\SI{5.0}{rad.s^{-1}}\). Find the first maximum transfer time.}{%
\fivechoices{\(\SI{1.57}{s}\)}{\(\SI{3.14}{s}\)}{\(\SI{4.00}{s}\)}{\(\SI{5.00}{s}\)}{\(\SI{6.28}{s}\)}}{%
\(t_{\mathrm{transfer}}=\pi/\Delta\omega=\pi/(5-4)=\SI{3.14}{s}\).}

\bankq{calc-mode-028}{Coupled SHM and Normal Modes}{challenge}{C}{generalized-eigenvalue,modes}{Calculate the lower generalized normal frequency.}{%
For \(M=2I\) and \(K=\begin{pmatrix}10&-4\\-4&10\end{pmatrix}\), find the lower normal angular frequency.}{%
\fivechoices{\(\SI{1.22}{rad.s^{-1}}\)}{\(\SI{1.41}{rad.s^{-1}}\)}{\(\SI{1.73}{rad.s^{-1}}\)}{\(\SI{2.24}{rad.s^{-1}}\)}{\(\SI{2.65}{rad.s^{-1}}\)}}{%
The stiffness eigenvalues are \(6\) and \(14\). Dividing by mass \(2\) gives \(\omega^2=3,7\), so \(\omega_{\min}=\sqrt3=\SI{1.73}{rad.s^{-1}}\).}

\bankq{calc-mode-029}{Coupled SHM and Normal Modes}{challenge}{D}{generalized-eigenvalue,modes}{Solve an unequal-mass normal-mode determinant.}{%
Let \(M=\operatorname{diag}(1,2)\) and \(K=\begin{pmatrix}3&-1\\-1&3\end{pmatrix}\) in SI units. Find the higher normal angular frequency.}{%
\fivechoices{\(\SI{1.10}{rad.s^{-1}}\)}{\(\SI{1.41}{rad.s^{-1}}\)}{\(\SI{1.58}{rad.s^{-1}}\)}{\(\SI{1.81}{rad.s^{-1}}\)}{\(\SI{2.28}{rad.s^{-1}}\)}}{%
\(\det(K-sM)=2s^2-9s+8=0\), so \(s=(9+\sqrt{17})/4=3.281\). Thus \(\omega=\sqrt{s}=\SI{1.81}{rad.s^{-1}}\).}

\bankq{calc-mode-030}{Coupled SHM and Normal Modes}{challenge}{E}{normal-modes,reconstruction}{Reconstruct a coordinate from normal coordinates.}{%
Normal coordinates are \(q_+=2\cos(3t)\) cm and \(q_-=4\cos(5t)\) cm, with \(x_1=(q_++q_-)/2\). Find \(x_1\) at \(t=\pi/10\) s.}{%
\fivechoices{\(-1.18\) cm}{\(-0.588\) cm}{0}{\(0.294\) cm}{\(0.588\) cm}}{%
At \(t=\pi/10\), \(q_+=2\cos(3\pi/10)=1.176\) cm and \(q_-=4\cos(\pi/2)=0\). Hence \(x_1=1.176/2=0.588\) cm.}

% 31--40: Travelling waves, power, and intensity (3 foundation, 5 applied, 2 challenge)
\bankq{calc-wave-031}{Travelling Waves, Power, and Intensity}{foundation}{A}{waves,speed}{Calculate wave speed.}{%
A wave has frequency \(\SI{25}{Hz}\) and wavelength \(\SI{0.80}{m}\). Find its speed.}{%
\fivechoices{\(\SI{20}{m.s^{-1}}\)}{\(\SI{25}{m.s^{-1}}\)}{\(\SI{31.3}{m.s^{-1}}\)}{\(\SI{40}{m.s^{-1}}\)}{\(\SI{50}{m.s^{-1}}\)}}{%
\(v=f\lambda=(25)(0.80)=\SI{20}{m.s^{-1}}\).}

\bankq{calc-wave-032}{Travelling Waves, Power, and Intensity}{foundation}{B}{waves,wavenumber}{Calculate wave number.}{%
Find the wave number of a sinusoidal wave with wavelength \(\SI{0.400}{m}\).}{%
\fivechoices{\(\SI{7.85}{m^{-1}}\)}{\(\SI{15.7}{m^{-1}}\)}{\(\SI{20.0}{m^{-1}}\)}{\(\SI{31.4}{m^{-1}}\)}{\(\SI{40.0}{m^{-1}}\)}}{%
\(k=2\pi/\lambda=2\pi/0.400=\SI{15.7}{m^{-1}}\).}

\bankq{calc-wave-033}{Travelling Waves, Power, and Intensity}{foundation}{C}{waves,angular-frequency}{Calculate angular frequency.}{%
A wave has frequency \(\SI{50.0}{Hz}\). Find its angular frequency.}{%
\fivechoices{\(\SI{50}{rad.s^{-1}}\)}{\(\SI{157}{rad.s^{-1}}\)}{\(\SI{314}{rad.s^{-1}}\)}{\(\SI{628}{rad.s^{-1}}\)}{\(\SI{2500}{rad.s^{-1}}\)}}{%
\(\omega=2\pi f=2\pi(50.0)=\SI{314}{rad.s^{-1}}\).}

\bankq{calc-wave-034}{Travelling Waves, Power, and Intensity}{applied}{D}{waves,phase-velocity}{Calculate signed phase velocity.}{%
For \(y=0.02\cos(6x+18t)\), with SI units, find the signed phase velocity along the \(x\)-axis.}{%
\fivechoices{\(\SI{9}{m.s^{-1}}\)}{\(\SI{6}{m.s^{-1}}\)}{\(\SI{3}{m.s^{-1}}\)}{\(\SI{-3}{m.s^{-1}}\)}{\(\SI{-18}{m.s^{-1}}\)}}{%
Constant phase gives \(6x+18t=\text{constant}\), hence \(dx/dt=-18/6=\SI{-3}{m.s^{-1}}\).}

\bankq{calc-wave-035}{Travelling Waves, Power, and Intensity}{applied}{E}{waves,particle-speed}{Calculate maximum particle speed.}{%
A transverse wave has amplitude \(\SI{0.020}{m}\) and angular frequency \(\SI{40}{rad.s^{-1}}\). Find the maximum particle speed.}{%
\fivechoices{\(\SI{0.20}{m.s^{-1}}\)}{\(\SI{0.40}{m.s^{-1}}\)}{\(\SI{0.60}{m.s^{-1}}\)}{\(\SI{0.70}{m.s^{-1}}\)}{\(\SI{0.80}{m.s^{-1}}\)}}{%
\(v_{y,\max}=A\omega=(0.020)(40)=\SI{0.80}{m.s^{-1}}\).}

\bankq{calc-wave-036}{Travelling Waves, Power, and Intensity}{applied}{A}{string,wave-speed}{Calculate string-wave speed.}{%
A string has tension \(\SI{144}{N}\) and linear density \(\SI{0.0100}{kg.m^{-1}}\). Find the wave speed.}{%
\fivechoices{\(\SI{120}{m.s^{-1}}\)}{\(\SI{100}{m.s^{-1}}\)}{\(\SI{72}{m.s^{-1}}\)}{\(\SI{14.4}{m.s^{-1}}\)}{\(\SI{12}{m.s^{-1}}\)}}{%
\(v=\sqrt{T/\mu}=\sqrt{144/0.0100}=\SI{120}{m.s^{-1}}\).}

\bankq{calc-wave-037}{Travelling Waves, Power, and Intensity}{applied}{B}{string,power}{Calculate average power on a string.}{%
A sinusoidal string wave has \(\mu=\SI{0.020}{kg.m^{-1}}\), \(\omega=\SI{50}{rad.s^{-1}}\), \(A=\SI{0.010}{m}\), and \(v=\SI{100}{m.s^{-1}}\). Find its average power.}{%
\fivechoices{\(\SI{0.125}{W}\)}{\(\SI{0.250}{W}\)}{\(\SI{0.500}{W}\)}{\(\SI{1.25}{W}\)}{\(\SI{2.50}{W}\)}}{%
\(\overline P=\tfrac12\mu\omega^2A^2v=\tfrac12(0.020)(50)^2(0.010)^2(100)=\SI{0.250}{W}\).}

\bankq{calc-wave-038}{Travelling Waves, Power, and Intensity}{applied}{C}{intensity,spherical-wave}{Calculate spherical-wave intensity.}{%
An isotropic source radiates \(\SI{12.6}{W}\). Find the intensity at \(r=\SI{10.0}{m}\).}{%
\fivechoices{\(\SI{1.00e-3}{W.m^{-2}}\)}{\(\SI{5.00e-3}{W.m^{-2}}\)}{\(\SI{1.00e-2}{W.m^{-2}}\)}{\(\SI{2.00e-2}{W.m^{-2}}\)}{\(\SI{1.00e-1}{W.m^{-2}}\)}}{%
\(I=P/(4\pi r^2)=12.6/[4\pi(10.0)^2]=\SI{1.00e-2}{W.m^{-2}}\).}

\bankq{calc-wave-039}{Travelling Waves, Power, and Intensity}{challenge}{D}{sound,pressure-amplitude}{Calculate pressure amplitude from intensity.}{%
A plane sound wave in air has intensity \(\SI{0.600}{W.m^{-2}}\), with \(\rho=\SI{1.20}{kg.m^{-3}}\) and \(v=\SI{343}{m.s^{-1}}\). Find its pressure amplitude.}{%
\fivechoices{\(\SI{11.1}{Pa}\)}{\(\SI{15.7}{Pa}\)}{\(\SI{19.8}{Pa}\)}{\(\SI{22.2}{Pa}\)}{\(\SI{31.4}{Pa}\)}}{%
\(p_0=\sqrt{2\rho vI}=\sqrt{2(1.20)(343)(0.600)}=\SI{22.2}{Pa}\).}

\bankq{calc-wave-040}{Travelling Waves, Power, and Intensity}{challenge}{E}{standing-waves,wavenumber}{Calculate wavelength and wave number from node spacing.}{%
Adjacent nodes of a standing wave are \(\SI{0.150}{m}\) apart. Find the wave number.}{%
\fivechoices{\(\SI{5.24}{m^{-1}}\)}{\(\SI{10.5}{m^{-1}}\)}{\(\SI{15.7}{m^{-1}}\)}{\(\SI{18.8}{m^{-1}}\)}{\(\SI{20.9}{m^{-1}}\)}}{%
Node spacing is \(\lambda/2\), so \(\lambda=\SI{0.300}{m}\) and \(k=2\pi/\lambda=\SI{20.9}{m^{-1}}\).}

% 41--50: Interference and beats (3 foundation, 4 applied, 3 challenge)
\bankq{calc-int-041}{Interference and Beats}{foundation}{A}{interference,resultant-amplitude}{Calculate resultant amplitude.}{%
Two waves each have amplitude \(\SI{3.0}{cm}\) and phase difference \(60^\circ\). Find the resultant amplitude.}{%
\fivechoices{\(\SI{5.20}{cm}\)}{\(\SI{4.24}{cm}\)}{\(\SI{3.00}{cm}\)}{\(\SI{2.60}{cm}\)}{\(\SI{1.50}{cm}\)}}{%
\(A_R=2A\cos(\phi/2)=2(3.0)\cos30^\circ=\SI{5.20}{cm}\).}

\bankq{calc-int-042}{Interference and Beats}{foundation}{B}{interference,phase}{Calculate phase difference from path difference.}{%
Two coherent waves have path difference \(0.75\lambda\). Find their phase difference in degrees.}{%
\fivechoices{\(180^\circ\)}{\(270^\circ\)}{\(300^\circ\)}{\(360^\circ\)}{\(540^\circ\)}}{%
\(\Delta\phi=2\pi\Delta x/\lambda=2\pi(0.75)=3\pi/2=270^\circ\).}

\bankq{calc-int-043}{Interference and Beats}{foundation}{C}{double-slit,fringe-spacing}{Calculate double-slit fringe spacing.}{%
Light of wavelength \(\SI{600}{nm}\) passes through slits separated by \(\SI{0.500}{mm}\). A screen is \(\SI{2.00}{m}\) away. Find the fringe spacing.}{%
\fivechoices{\(\SI{0.60}{mm}\)}{\(\SI{1.20}{mm}\)}{\(\SI{2.40}{mm}\)}{\(\SI{4.80}{mm}\)}{\(\SI{6.00}{mm}\)}}{%
\(\Delta y=\lambda L/d=(600\times10^{-9})(2.00)/(0.500\times10^{-3})=\SI{2.40}{mm}\).}

\bankq{calc-int-044}{Interference and Beats}{applied}{D}{double-slit,bright-fringe}{Calculate bright-fringe position.}{%
Using the arrangement in Question 43, find the position of the fourth bright fringe from the centre.}{%
\fivechoices{\(\SI{2.4}{mm}\)}{\(\SI{4.8}{mm}\)}{\(\SI{7.2}{mm}\)}{\(\SI{9.6}{mm}\)}{\(\SI{12.0}{mm}\)}}{%
\(y_4=4\Delta y=4(2.40)=\SI{9.60}{mm}\).}

\bankq{calc-int-045}{Interference and Beats}{applied}{E}{beats,frequency}{Calculate beat frequency.}{%
Two tuning forks have frequencies \(\SI{256}{Hz}\) and \(\SI{260}{Hz}\). Find the beat frequency.}{%
\fivechoices{\(\SI{516}{Hz}\)}{\(\SI{258}{Hz}\)}{\(\SI{8}{Hz}\)}{\(\SI{6}{Hz}\)}{\(\SI{4}{Hz}\)}}{%
\(f_b=|f_2-f_1|=|260-256|=\SI{4}{Hz}\).}

\bankq{calc-int-046}{Interference and Beats}{applied}{A}{interference,intensity}{Calculate constructive intensity.}{%
Two coherent waves separately produce intensities \(\SI{9}{W.m^{-2}}\) and \(\SI{4}{W.m^{-2}}\). Find the intensity at complete constructive interference.}{%
\fivechoices{\(\SI{25}{W.m^{-2}}\)}{\(\SI{13}{W.m^{-2}}\)}{\(\SI{10}{W.m^{-2}}\)}{\(\SI{5}{W.m^{-2}}\)}{\(\SI{1}{W.m^{-2}}\)}}{%
Amplitudes are proportional to \(\sqrt I\), so \(I_{\max}=(\sqrt9+\sqrt4)^2=(3+2)^2=\SI{25}{W.m^{-2}}\).}

\bankq{calc-int-047}{Interference and Beats}{applied}{B}{interference,intensity}{Calculate destructive intensity.}{%
For the same two waves as Question 46, find the intensity at complete destructive interference.}{%
\fivechoices{\(\SI{0}{W.m^{-2}}\)}{\(\SI{1}{W.m^{-2}}\)}{\(\SI{5}{W.m^{-2}}\)}{\(\SI{9}{W.m^{-2}}\)}{\(\SI{13}{W.m^{-2}}\)}}{%
\(I_{\min}=(\sqrt9-\sqrt4)^2=(3-2)^2=\SI{1}{W.m^{-2}}\).}

\bankq{calc-int-048}{Interference and Beats}{challenge}{C}{interference,phase}{Infer phase difference from intensity.}{%
Two equal coherent sources each produce intensity \(I_0\). Their resultant intensity is \(I_0\). Find the smaller phase difference.}{%
\fivechoices{\(30^\circ\)}{\(60^\circ\)}{\(120^\circ\)}{\(150^\circ\)}{\(180^\circ\)}}{%
\(I=2I_0(1+\cos\phi)=I_0\), so \(\cos\phi=-1/2\) and \(\phi=120^\circ\).}

\bankq{calc-int-049}{Interference and Beats}{challenge}{D}{thin-film,interference}{Calculate minimum film thickness.}{%
A film of refractive index \(1.50\) is in air and illuminated normally with \(\SI{600}{nm}\) light. With one phase reversal, find the least nonzero thickness for constructive reflected light.}{%
\fivechoices{\(\SI{50}{nm}\)}{\(\SI{75}{nm}\)}{\(\SI{80}{nm}\)}{\(\SI{100}{nm}\)}{\(\SI{200}{nm}\)}}{%
With one phase reversal, the first constructive condition is \(2nt=\lambda/2\). Thus \(t=\lambda/(4n)=600/[4(1.50)]=\SI{100}{nm}\).}

\bankq{calc-int-050}{Interference and Beats}{challenge}{E}{interferometer,fringes}{Calculate Michelson fringe count.}{%
In a Michelson interferometer using \(\SI{600}{nm}\) light, one mirror moves \(\SI{0.150}{mm}\). How many fringes pass the detector?}{%
\fivechoices{125}{250}{300}{400}{500}{}}{%
\(N=2\Delta d/\lambda=2(0.150\times10^{-3})/(600\times10^{-9})=500\).}

% 51--60: Diffraction and polarization (3 foundation, 5 applied, 2 challenge)
\bankq{calc-diff-051}{Diffraction and Polarization}{foundation}{A}{single-slit,minima}{Calculate a single-slit minimum angle.}{%
Light of wavelength \(\SI{500}{nm}\) passes through a slit of width \(\SI{0.100}{mm}\). Find the first-minimum angle.}{%
\fivechoices{\(0.286^\circ\)}{\(0.410^\circ\)}{\(0.500^\circ\)}{\(0.573^\circ\)}{\(1.00^\circ\)}}{%
\(\sin\theta_1=\lambda/a=500\times10^{-9}/10^{-4}=0.005\), so \(\theta_1=0.286^\circ\).}

\bankq{calc-diff-052}{Diffraction and Polarization}{foundation}{B}{single-slit,central-width}{Calculate central diffraction width.}{%
A slit of width \(\SI{0.200}{mm}\) is illuminated by \(\SI{600}{nm}\) light. A screen is \(\SI{2.00}{m}\) away. Find the central-maximum width.}{%
\fivechoices{\(\SI{6.0}{mm}\)}{\(\SI{12.0}{mm}\)}{\(\SI{18.0}{mm}\)}{\(\SI{24.0}{mm}\)}{\(\SI{30.0}{mm}\)}}{%
For small angles, \(w=2L\lambda/a=2(2.00)(600\times10^{-9})/(0.200\times10^{-3})=\SI{12.0}{mm}\).}

\bankq{calc-diff-053}{Diffraction and Polarization}{foundation}{C}{single-slit,minima}{Calculate a higher diffraction minimum.}{%
For the slit in Question 51, find the angle of the second minimum.}{%
\fivechoices{\(0.143^\circ\)}{\(0.286^\circ\)}{\(0.573^\circ\)}{\(0.860^\circ\)}{\(1.15^\circ\)}}{%
\(\sin\theta_2=2\lambda/a=0.010\), giving \(\theta_2=0.573^\circ\).}

\bankq{calc-diff-054}{Diffraction and Polarization}{applied}{D}{single-slit,maxima}{Calculate a secondary diffraction maximum.}{%
For \(\lambda=\SI{500}{nm}\) and slit width \(a=\SI{0.100}{mm}\), use the first root \(\beta=4.493\) of \(\tan\beta=\beta\). Find the first secondary-maximum angle.}{%
\fivechoices{\(0.143^\circ\)}{\(0.286^\circ\)}{\(0.358^\circ\)}{\(0.410^\circ\)}{\(0.573^\circ\)}}{%
Since \(\beta=\pi a\sin\theta/\lambda\), \(\sin\theta=4.493(500\times10^{-9})/[\pi(10^{-4})]=0.00715\). Thus \(\theta=0.410^\circ\).}

\bankq{calc-diff-055}{Diffraction and Polarization}{applied}{E}{grating,order}{Calculate the greatest grating order.}{%
A grating has spacing \(d=\SI{2.00}{um}\) and is illuminated by \(\lambda=\SI{500}{nm}\). What is the greatest possible order?}{%
\fivechoices{0}{1}{2}{3}{4}{}}{%
Because \(m\lambda\le d\), \(m_{\max}=\lfloor d/\lambda\rfloor=\lfloor2000/500\rfloor=4\).}

\bankq{calc-diff-056}{Diffraction and Polarization}{applied}{A}{grating,angle}{Calculate a grating diffraction angle.}{%
A grating has \(1000\) lines per mm and is illuminated normally by \(\SI{600}{nm}\) light. Find the first-order angle.}{%
\fivechoices{\(36.9^\circ\)}{\(30.0^\circ\)}{\(26.6^\circ\)}{\(18.4^\circ\)}{\(12.0^\circ\)}}{%
The spacing is \(d=1/(10^6)=\SI{1.00}{um}\). Thus \(\sin\theta=\lambda/d=0.600\), so \(\theta=36.9^\circ\).}

\bankq{calc-diff-057}{Diffraction and Polarization}{applied}{B}{circular-aperture,resolution}{Calculate Rayleigh angular resolution.}{%
A circular aperture has diameter \(\SI{0.100}{m}\) and observes \(\SI{500}{nm}\) light. Find its Rayleigh angular resolution.}{%
\fivechoices{\(\SI{5.00e-6}{rad}\)}{\(\SI{6.10e-6}{rad}\)}{\(\SI{1.00e-5}{rad}\)}{\(\SI{1.22e-5}{rad}\)}{\(\SI{6.10e-5}{rad}\)}}{%
\(\theta_R=1.22\lambda/D=1.22(500\times10^{-9})/0.100=\SI{6.10e-6}{rad}\).}

\bankq{calc-diff-058}{Diffraction and Polarization}{applied}{C}{polarization,malus-law}{Calculate transmitted polarized intensity.}{%
Plane-polarized light of intensity \(\SI{80}{W.m^{-2}}\) meets an analyzer at \(30^\circ\). Find the transmitted intensity.}{%
\fivechoices{\(\SI{20}{W.m^{-2}}\)}{\(\SI{40}{W.m^{-2}}\)}{\(\SI{60}{W.m^{-2}}\)}{\(\SI{69.3}{W.m^{-2}}\)}{\(\SI{80}{W.m^{-2}}\)}}{%
By Malus's law, \(I=I_0\cos^230^\circ=80(3/4)=\SI{60}{W.m^{-2}}\).}

\bankq{calc-diff-059}{Diffraction and Polarization}{challenge}{D}{polarization,malus-law}{Calculate transmission through two polarizers.}{%
Unpolarized light passes through an ideal polarizer and then an analyzer at \(45^\circ\). What percentage of the original intensity emerges?}{%
\fivechoices{12.5\%}{18.8\%}{20.0\%}{25.0\%}{50.0\%}}{%
The first polarizer gives \(I_0/2\), and the analyzer gives another factor \(\cos^245^\circ=1/2\). Thus \(I=I_0/4=25.0\%\,I_0\).}

\bankq{calc-diff-060}{Diffraction and Polarization}{challenge}{E}{polarization,brewster-angle}{Calculate Brewster angle.}{%
Light travels from air onto glass of refractive index \(1.50\). Find Brewster's angle.}{%
\fivechoices{\(33.7^\circ\)}{\(41.8^\circ\)}{\(45.0^\circ\)}{\(48.2^\circ\)}{\(56.3^\circ\)}}{%
\(\tan\theta_B=n_2/n_1=1.50\), so \(\theta_B=\tan^{-1}(1.50)=56.3^\circ\).}

% 61--70: Dispersion and velocities (3 foundation, 4 applied, 3 challenge)
\bankq{calc-disp-061}{Dispersion, Phase Velocity, and Group Velocity}{foundation}{A}{dispersion,phase-velocity}{Calculate phase velocity from a dispersion relation.}{%
A medium obeys \(\omega=4k\) in SI units. Find the phase velocity.}{%
\fivechoices{\(\SI{4.0}{m.s^{-1}}\)}{\(\SI{2.0}{m.s^{-1}}\)}{\(\SI{1.0}{m.s^{-1}}\)}{\(\SI{0.50}{m.s^{-1}}\)}{\(\SI{0.25}{m.s^{-1}}\)}}{%
\(v_p=\omega/k=4k/k=\SI{4.0}{m.s^{-1}}\).}

\bankq{calc-disp-062}{Dispersion, Phase Velocity, and Group Velocity}{foundation}{B}{dispersion,group-velocity}{Calculate group velocity by differentiation.}{%
For \(\omega=2k+0.5k^2\), find the group velocity at \(k=\SI{3.0}{m^{-1}}\).}{%
\fivechoices{\(\SI{3.5}{m.s^{-1}}\)}{\(\SI{5.0}{m.s^{-1}}\)}{\(\SI{6.0}{m.s^{-1}}\)}{\(\SI{7.5}{m.s^{-1}}\)}{\(\SI{10.5}{m.s^{-1}}\)}}{%
\(v_g=d\omega/dk=2+k\). At \(k=3.0\), \(v_g=\SI{5.0}{m.s^{-1}}\).}

\bankq{calc-disp-063}{Dispersion, Phase Velocity, and Group Velocity}{foundation}{C}{dispersion,inverse-relation}{Calculate group velocity from k as a function of omega.}{%
A medium obeys \(k=0.040\omega^2\). Find \(v_g\) at \(\omega=\SI{10}{rad.s^{-1}}\).}{%
\fivechoices{\(\SI{0.80}{m.s^{-1}}\)}{\(\SI{1.00}{m.s^{-1}}\)}{\(\SI{1.25}{m.s^{-1}}\)}{\(\SI{2.50}{m.s^{-1}}\)}{\(\SI{4.00}{m.s^{-1}}\)}}{%
\(dk/d\omega=0.080\omega=0.80\). Therefore \(v_g=d\omega/dk=(dk/d\omega)^{-1}=\SI{1.25}{m.s^{-1}}\).}

\bankq{calc-disp-064}{Dispersion, Phase Velocity, and Group Velocity}{applied}{D}{water-waves,phase-velocity}{Calculate deep-water phase velocity.}{%
Deep-water waves obey \(\omega=\sqrt{gk}\). Find \(v_p\) at \(k=\SI{4.00}{m^{-1}}\).}{%
\fivechoices{\(\SI{0.783}{m.s^{-1}}\)}{\(\SI{1.00}{m.s^{-1}}\)}{\(\SI{1.25}{m.s^{-1}}\)}{\(\SI{1.57}{m.s^{-1}}\)}{\(\SI{3.13}{m.s^{-1}}\)}}{%
\(\omega=\sqrt{9.80(4.00)}=\SI{6.26}{rad.s^{-1}}\), so \(v_p=\omega/k=6.26/4.00=\SI{1.57}{m.s^{-1}}\).}

\bankq{calc-disp-065}{Dispersion, Phase Velocity, and Group Velocity}{applied}{E}{water-waves,group-velocity}{Calculate deep-water group velocity.}{%
For the wave in Question 64, find the group velocity.}{%
\fivechoices{\(\SI{3.13}{m.s^{-1}}\)}{\(\SI{1.57}{m.s^{-1}}\)}{\(\SI{1.25}{m.s^{-1}}\)}{\(\SI{1.00}{m.s^{-1}}\)}{\(\SI{0.783}{m.s^{-1}}\)}}{%
For \(\omega=\sqrt{gk}\), \(v_g=\tfrac12v_p=\tfrac12(1.57)=\SI{0.783}{m.s^{-1}}\).}

\bankq{calc-disp-066}{Dispersion, Phase Velocity, and Group Velocity}{applied}{A}{dispersion,cutoff}{Calculate group velocity with a cutoff frequency.}{%
A medium obeys \(\omega^2=9k^2+25\). Find \(v_g\) at \(k=\SI{4.0}{m^{-1}}\).}{%
\fivechoices{\(\SI{2.77}{m.s^{-1}}\)}{\(\SI{3.00}{m.s^{-1}}\)}{\(\SI{3.25}{m.s^{-1}}\)}{\(\SI{4.00}{m.s^{-1}}\)}{\(\SI{13.0}{m.s^{-1}}\)}}{%
Here \(\omega=\sqrt{9(16)+25}=13\). Differentiation gives \(v_g=9k/\omega=36/13=\SI{2.77}{m.s^{-1}}\).}

\bankq{calc-disp-067}{Dispersion, Phase Velocity, and Group Velocity}{applied}{B}{optics,phase-velocity}{Calculate optical phase speed.}{%
Light travels in a medium of refractive index \(1.50\). Find its phase speed.}{%
\fivechoices{\(\SI{1.50e8}{m.s^{-1}}\)}{\(\SI{2.00e8}{m.s^{-1}}\)}{\(\SI{2.25e8}{m.s^{-1}}\)}{\(\SI{3.00e8}{m.s^{-1}}\)}{\(\SI{4.50e8}{m.s^{-1}}\)}}{%
\(v_p=c/n=(3.00\times10^8)/1.50=\SI{2.00e8}{m.s^{-1}}\).}

\bankq{calc-disp-068}{Dispersion, Phase Velocity, and Group Velocity}{challenge}{C}{dispersion,velocity-ratio}{Calculate the group-to-phase velocity ratio.}{%
A medium obeys \(\omega=0.20k^3\). Find \(v_g/v_p\) at \(k=\SI{2.0}{m^{-1}}\).}{%
\fivechoices{1}{2}{3}{4}{6}{}}{%
\(v_p=\omega/k=0.20k^2\), while \(v_g=d\omega/dk=0.60k^2\). Thus \(v_g/v_p=3\).}

\bankq{calc-disp-069}{Dispersion, Phase Velocity, and Group Velocity}{challenge}{D}{dispersion,pulse-spreading}{Estimate dispersive pulse spreading.}{%
For \(\omega=ck+0.50k^2\), a packet has wave-number width \(\Delta k=\SI{0.40}{m^{-1}}\). Estimate its extra width after \(\SI{10}{s}\) using \(\Delta v_g\approx|d^2\omega/dk^2|\Delta k\).}{%
\fivechoices{\(\SI{1.0}{m}\)}{\(\SI{2.0}{m}\)}{\(\SI{3.0}{m}\)}{\(\SI{4.0}{m}\)}{\(\SI{5.0}{m}\)}}{%
\(d^2\omega/dk^2=1.0\), so \(\Delta v_g=(1.0)(0.40)=\SI{0.40}{m.s^{-1}}\). Hence \(\Delta x\approx\Delta v_gt=0.40(10)=\SI{4.0}{m}\).}

\bankq{calc-disp-070}{Dispersion, Phase Velocity, and Group Velocity}{challenge}{E}{dispersion,implicit-differentiation}{Calculate group velocity by implicit differentiation.}{%
A dispersion relation is \(\omega^2-9k^2-0.10k^4=0\). Find \(v_g\) at \(k=\SI{2.0}{m^{-1}}\).}{%
\fivechoices{\(\SI{2.40}{m.s^{-1}}\)}{\(\SI{2.77}{m.s^{-1}}\)}{\(\SI{3.00}{m.s^{-1}}\)}{\(\SI{3.13}{m.s^{-1}}\)}{\(\SI{3.20}{m.s^{-1}}\)}}{%
At \(k=2\), \(\omega=\sqrt{36+1.6}=6.132\). Differentiation gives \(2\omega v_g=18k+0.40k^3\), so \(v_g=(36+3.2)/(2\times6.132)=\SI{3.20}{m.s^{-1}}\).}

% 71--80: Wave equations and bounded modes (3 foundation, 5 applied, 2 challenge)
\bankq{calc-eq-071}{Wave Equations and Bounded Modes}{foundation}{A}{wave-equation,speed}{Calculate speed from a 1D wave equation.}{%
A field obeys \(u_{tt}=225u_{xx}\). Find the wave speed.}{%
\fivechoices{\(\SI{15}{m.s^{-1}}\)}{\(\SI{112.5}{m.s^{-1}}\)}{\(\SI{225}{m.s^{-1}}\)}{\(\SI{450}{m.s^{-1}}\)}{\(\SI{50625}{m.s^{-1}}\)}}{%
Comparing with \(u_{tt}=v^2u_{xx}\) gives \(v=\sqrt{225}=\SI{15}{m.s^{-1}}\).}

\bankq{calc-eq-072}{Wave Equations and Bounded Modes}{foundation}{B}{string,fundamental}{Calculate a string fundamental frequency.}{%
A string of length \(\SI{1.20}{m}\) is fixed at both ends and supports waves at \(\SI{120}{m.s^{-1}}\). Find its fundamental frequency.}{%
\fivechoices{\(\SI{25}{Hz}\)}{\(\SI{50}{Hz}\)}{\(\SI{75}{Hz}\)}{\(\SI{100}{Hz}\)}{\(\SI{144}{Hz}\)}}{%
\(f_1=v/(2L)=120/[2(1.20)]=\SI{50}{Hz}\).}

\bankq{calc-eq-073}{Wave Equations and Bounded Modes}{foundation}{C}{string,harmonics}{Calculate a string harmonic.}{%
For the string in Question 72, find the third-harmonic frequency.}{%
\fivechoices{\(\SI{50}{Hz}\)}{\(\SI{100}{Hz}\)}{\(\SI{150}{Hz}\)}{\(\SI{200}{Hz}\)}{\(\SI{300}{Hz}\)}}{%
\(f_3=3f_1=3(50)=\SI{150}{Hz}\).}

\bankq{calc-eq-074}{Wave Equations and Bounded Modes}{applied}{D}{membrane,normal-mode}{Calculate a rectangular membrane frequency.}{%
A rectangular membrane has sides \(a=\SI{2.0}{m}\), \(b=\SI{1.0}{m}\), and wave speed \(v=\SI{40}{m.s^{-1}}\). Find the \((1,1)\) mode frequency.}{%
\fivechoices{\(\SI{10.0}{Hz}\)}{\(\SI{15.0}{Hz}\)}{\(\SI{20.0}{Hz}\)}{\(\SI{22.4}{Hz}\)}{\(\SI{28.3}{Hz}\)}}{%
\(f_{mn}=\tfrac v2\sqrt{(m/a)^2+(n/b)^2}\). Thus \(f_{11}=20\sqrt{(1/2)^2+1^2}=\SI{22.4}{Hz}\).}

\bankq{calc-eq-075}{Wave Equations and Bounded Modes}{applied}{E}{membrane,normal-mode}{Calculate a second membrane mode.}{%
For the membrane in Question 74, find the \((2,1)\) mode frequency.}{%
\fivechoices{\(\SI{20.0}{Hz}\)}{\(\SI{22.4}{Hz}\)}{\(\SI{25.0}{Hz}\)}{\(\SI{26.8}{Hz}\)}{\(\SI{28.3}{Hz}\)}}{%
\(f_{21}=20\sqrt{(2/2)^2+(1/1)^2}=20\sqrt2=\SI{28.3}{Hz}\).}

\bankq{calc-eq-076}{Wave Equations and Bounded Modes}{applied}{A}{cavity,3d-mode}{Calculate a cubic-cavity mode frequency.}{%
A cubic cavity of side \(\SI{2.0}{m}\) supports waves at \(\SI{300}{m.s^{-1}}\). Find the \((1,1,1)\) mode frequency using fixed-boundary mode numbers.}{%
\fivechoices{\(\SI{130}{Hz}\)}{\(\SI{150}{Hz}\)}{\(\SI{212}{Hz}\)}{\(\SI{260}{Hz}\)}{\(\SI{450}{Hz}\)}}{%
\(f_{111}=\tfrac v2\sqrt{3(1/2)^2}=150(\sqrt3/2)=\SI{130}{Hz}\).}

\bankq{calc-eq-077}{Wave Equations and Bounded Modes}{applied}{B}{wave-vector,wavelength}{Calculate wavelength from a 3D wave vector.}{%
A plane wave has \(\mathbf{k}=(3,4,12)\,\si{m^{-1}}\). Find its wavelength.}{%
\fivechoices{\(\SI{0.242}{m}\)}{\(\SI{0.483}{m}\)}{\(\SI{0.785}{m}\)}{\(\SI{1.21}{m}\)}{\(\SI{2.07}{m}\)}}{%
\(|\mathbf{k}|=\sqrt{3^2+4^2+12^2}=13\,\si{m^{-1}}\), so \(\lambda=2\pi/13=\SI{0.483}{m}\).}

\bankq{calc-eq-078}{Wave Equations and Bounded Modes}{applied}{C}{wave-equation,plane-wave}{Calculate wave speed from omega and wave vector.}{%
A plane-wave solution of \(u_{tt}=v^2\nabla^2u\) has \(|\mathbf{k}|=\SI{5.0}{m^{-1}}\) and \(\omega=\SI{50}{rad.s^{-1}}\). Find \(v\).}{%
\fivechoices{\(\SI{2.0}{m.s^{-1}}\)}{\(\SI{5.0}{m.s^{-1}}\)}{\(\SI{10.0}{m.s^{-1}}\)}{\(\SI{25.0}{m.s^{-1}}\)}{\(\SI{250}{m.s^{-1}}\)}}{%
For a nondispersive wave, \(\omega=v|\mathbf{k}|\). Thus \(v=50/5.0=\SI{10.0}{m.s^{-1}}\).}

\bankq{calc-eq-079}{Wave Equations and Bounded Modes}{challenge}{D}{dalembert,pulse}{Track a travelling pulse.}{%
The peak of \(f(x-6t)\) is at \(x=-\SI{2.0}{m}\) when \(t=0\). Where is it at \(t=\SI{3.0}{s}\)?}{%
\fivechoices{\(-\SI{20}{m}\)}{\(-\SI{16}{m}\)}{\(\SI{12}{m}\)}{\(\SI{16}{m}\)}{\(\SI{20}{m}\)}}{%
The pulse travels in \(+x\) at \(\SI{6}{m.s^{-1}}\), so \(x=-2+6(3)=\SI{16}{m}\).}

\bankq{calc-eq-080}{Wave Equations and Bounded Modes}{challenge}{E}{spherical-wave,travel-time}{Calculate 3D wave travel time.}{%
A spherical wave travels \(\SI{30}{m}\) through a medium at \(\SI{1500}{m.s^{-1}}\). Find its travel time.}{%
\fivechoices{\(\SI{2.0}{ms}\)}{\(\SI{5.0}{ms}\)}{\(\SI{10}{ms}\)}{\(\SI{15}{ms}\)}{\(\SI{20}{ms}\)}}{%
\(t=r/v=30/1500=0.020\,\si{s}=\SI{20}{ms}\).}

% 81--90: Acoustics (3 foundation, 4 applied, 3 challenge)
\bankq{calc-ac-081}{Acoustics}{foundation}{A}{echo,distance}{Calculate echo distance.}{%
An echo returns \(\SI{1.20}{s}\) after a shout. How far away is the reflector?}{%
\fivechoices{\(\SI{206}{m}\)}{\(\SI{286}{m}\)}{\(\SI{343}{m}\)}{\(\SI{412}{m}\)}{\(\SI{686}{m}\)}}{%
The sound makes a round trip, so \(d=vt/2=343(1.20)/2=\SI{206}{m}\).}

\bankq{calc-ac-082}{Acoustics}{foundation}{B}{beats,frequency}{Calculate an unknown beat frequency.}{%
A \(\SI{440}{Hz}\) fork produces \(\SI{6}{beats.s^{-1}}\) with a higher-frequency fork. Find the second frequency.}{%
\fivechoices{\(\SI{434}{Hz}\)}{\(\SI{446}{Hz}\)}{\(\SI{440}{Hz}\)}{\(\SI{452}{Hz}\)}{\(\SI{2640}{Hz}\)}}{%
Since the unknown fork is higher, \(f_2=440+6=\SI{446}{Hz}\).}

\bankq{calc-ac-083}{Acoustics}{foundation}{C}{doppler,moving-source}{Calculate Doppler shift for a moving source.}{%
A \(\SI{500}{Hz}\) source approaches a stationary observer at \(\SI{20}{m.s^{-1}}\). Find the observed frequency.}{%
\fivechoices{\(\SI{472}{Hz}\)}{\(\SI{500}{Hz}\)}{\(\SI{531}{Hz}\)}{\(\SI{563}{Hz}\)}{\(\SI{686}{Hz}\)}}{%
\(f'=fv/(v-v_s)=500(343)/(343-20)=\SI{531}{Hz}\).}

\bankq{calc-ac-084}{Acoustics}{applied}{D}{doppler,moving-observer}{Calculate Doppler shift for a moving observer.}{%
A stationary \(\SI{600}{Hz}\) source is approached by an observer at \(\SI{15}{m.s^{-1}}\). Find the observed frequency.}{%
\fivechoices{\(\SI{574}{Hz}\)}{\(\SI{600}{Hz}\)}{\(\SI{615}{Hz}\)}{\(\SI{626}{Hz}\)}{\(\SI{643}{Hz}\)}}{%
\(f'=f(v+v_o)/v=600(343+15)/343=\SI{626}{Hz}\).}

\bankq{calc-ac-085}{Acoustics}{applied}{E}{doppler,moving-source-observer}{Calculate Doppler shift with both moving.}{%
A \(\SI{400}{Hz}\) source moves toward an observer at \(\SI{20}{m.s^{-1}}\), while the observer moves toward the source at \(\SI{10}{m.s^{-1}}\). Find the observed frequency.}{%
\fivechoices{\(\SI{388}{Hz}\)}{\(\SI{400}{Hz}\)}{\(\SI{412}{Hz}\)}{\(\SI{425}{Hz}\)}{\(\SI{437}{Hz}\)}}{%
\(f'=f(v+v_o)/(v-v_s)=400(343+10)/(343-20)=\SI{437}{Hz}\).}

\bankq{calc-ac-086}{Acoustics}{applied}{A}{closed-pipe,fundamental}{Calculate a closed-pipe fundamental.}{%
A pipe closed at one end has length \(\SI{0.850}{m}\). Find its fundamental frequency.}{%
\fivechoices{\(\SI{101}{Hz}\)}{\(\SI{202}{Hz}\)}{\(\SI{303}{Hz}\)}{\(\SI{404}{Hz}\)}{\(\SI{505}{Hz}\)}}{%
\(f_1=v/(4L)=343/[4(0.850)]=\SI{101}{Hz}\).}

\bankq{calc-ac-087}{Acoustics}{applied}{B}{open-pipe,harmonics}{Calculate an open-pipe harmonic.}{%
An open pipe is \(\SI{1.00}{m}\) long. Find its third-harmonic frequency.}{%
\fivechoices{\(\SI{343}{Hz}\)}{\(\SI{515}{Hz}\)}{\(\SI{686}{Hz}\)}{\(\SI{858}{Hz}\)}{\(\SI{1029}{Hz}\)}}{%
For an open pipe, \(f_3=3v/(2L)=3(343)/2=\SI{515}{Hz}\).}

\bankq{calc-ac-088}{Acoustics}{challenge}{C}{sound-level,decibel}{Calculate sound intensity level.}{%
A sound has intensity \(\SI{1.00e-6}{W.m^{-2}}\). Find its intensity level.}{%
\fivechoices{\(\SI{20}{dB}\)}{\(\SI{40}{dB}\)}{\(\SI{60}{dB}\)}{\(\SI{80}{dB}\)}{\(\SI{120}{dB}\)}}{%
\(\beta=10\log_{10}(I/I_{\mathrm{ref}})=10\log_{10}(10^{-6}/10^{-12})=\SI{60}{dB}\).}

\bankq{calc-ac-089}{Acoustics}{challenge}{D}{sound-level,addition}{Add two sound levels.}{%
Two independent sources each produce \(\SI{70.0}{dB}\) at a point. Find the combined level.}{%
\fivechoices{\(\SI{70.0}{dB}\)}{\(\SI{71.0}{dB}\)}{\(\SI{72.0}{dB}\)}{\(\SI{73.0}{dB}\)}{\(\SI{140}{dB}\)}}{%
Equal independent intensities add, so the increase is \(10\log_{10}2=\SI{3.01}{dB}\). The total is \(70.0+3.01=\SI{73.0}{dB}\).}

\bankq{calc-ac-090}{Acoustics}{challenge}{E}{well,fall-time,sound-time}{Calculate depth using fall and sound times.}{%
A stone is dropped into a well and the splash is heard \(\SI{3.00}{s}\) later. Neglect air resistance. Find the depth by solving \(\sqrt{2h/g}+h/343=3.00\).}{%
\fivechoices{\(\SI{35.0}{m}\)}{\(\SI{37.8}{m}\)}{\(\SI{39.2}{m}\)}{\(\SI{40.0}{m}\)}{\(\SI{40.7}{m}\)}}{%
Testing \(h=\SI{40.7}{m}\) gives \(t_f=\sqrt{2(40.7)/9.80}=2.882\,\si{s}\) and \(t_s=40.7/343=0.119\,\si{s}\). Their sum is \(3.001\,\si{s}\), so \(h\approx\SI{40.7}{m}\).}

% 91--100: Propagation in different media (3 foundation, 5 applied, 2 challenge)
\bankq{calc-med-091}{Propagation in Different Media}{foundation}{A}{string,propagation}{Calculate transverse speed on a string.}{%
A string has tension \(\SI{196}{N}\) and linear density \(\SI{0.0140}{kg.m^{-1}}\). Find its wave speed.}{%
\fivechoices{\(\SI{118}{m.s^{-1}}\)}{\(\SI{98.0}{m.s^{-1}}\)}{\(\SI{84.0}{m.s^{-1}}\)}{\(\SI{14.0}{m.s^{-1}}\)}{\(\SI{3.74}{m.s^{-1}}\)}}{%
\(v=\sqrt{T/\mu}=\sqrt{196/0.0140}=\sqrt{14000}=\SI{118}{m.s^{-1}}\).}

\bankq{calc-med-092}{Propagation in Different Media}{foundation}{B}{liquid,sound-speed}{Calculate sound speed in a liquid.}{%
Water has bulk modulus \(B=\SI{2.25e9}{Pa}\) and density \(\rho=\SI{1000}{kg.m^{-3}}\). Find the sound speed.}{%
\fivechoices{\(\SI{1125}{m.s^{-1}}\)}{\(\SI{1500}{m.s^{-1}}\)}{\(\SI{2250}{m.s^{-1}}\)}{\(\SI{3000}{m.s^{-1}}\)}{\(\SI{4.74e4}{m.s^{-1}}\)}}{%
\(v=\sqrt{B/\rho}=\sqrt{2.25\times10^9/1000}=\SI{1500}{m.s^{-1}}\).}

\bankq{calc-med-093}{Propagation in Different Media}{foundation}{C}{solid,rod-wave}{Calculate longitudinal speed in a solid rod.}{%
A rod has Young's modulus \(Y=\SI{7.20e10}{Pa}\) and density \(\rho=\SI{8000}{kg.m^{-3}}\). Find the longitudinal wave speed.}{%
\fivechoices{\(\SI{1500}{m.s^{-1}}\)}{\(\SI{2400}{m.s^{-1}}\)}{\(\SI{3000}{m.s^{-1}}\)}{\(\SI{6000}{m.s^{-1}}\)}{\(\SI{9000}{m.s^{-1}}\)}}{%
\(v=\sqrt{Y/\rho}=\sqrt{7.20\times10^{10}/8000}=\sqrt{9.00\times10^6}=\SI{3000}{m.s^{-1}}\).}

\bankq{calc-med-094}{Propagation in Different Media}{applied}{D}{gas,sound-speed}{Calculate sound speed in a gas.}{%
For air, take \(\gamma=1.40\), \(P=\SI{1.01e5}{Pa}\), and \(\rho=\SI{1.20}{kg.m^{-3}}\). Find \(v=\sqrt{\gamma P/\rho}\).}{%
\fivechoices{\(\SI{289}{m.s^{-1}}\)}{\(\SI{316}{m.s^{-1}}\)}{\(\SI{331}{m.s^{-1}}\)}{\(\SI{343}{m.s^{-1}}\)}{\(\SI{412}{m.s^{-1}}\)}}{%
\(v=\sqrt{1.40(1.01\times10^5)/1.20}=\SI{343}{m.s^{-1}}\).}

\bankq{calc-med-095}{Propagation in Different Media}{applied}{E}{impedance,acoustics}{Calculate acoustic impedance.}{%
Water has density \(\SI{1000}{kg.m^{-3}}\) and sound speed \(\SI{1500}{m.s^{-1}}\). Find its acoustic impedance.}{%
\fivechoices{\(\SI{1.50e3}{kg.m^{-2}.s^{-1}}\)}{\(\SI{1.50e4}{kg.m^{-2}.s^{-1}}\)}{\(\SI{1.50e5}{kg.m^{-2}.s^{-1}}\)}{\(\SI{1.00e6}{kg.m^{-2}.s^{-1}}\)}{\(\SI{1.50e6}{kg.m^{-2}.s^{-1}}\)}}{%
\(Z=\rho v=(1000)(1500)=\SI{1.50e6}{kg.m^{-2}.s^{-1}}\).}

\bankq{calc-med-096}{Propagation in Different Media}{applied}{A}{impedance,reflection}{Calculate pressure reflection coefficient.}{%
A wave meets a boundary from impedance \(Z_1=2.0\) units to \(Z_2=6.0\) units. Find the pressure-amplitude reflection coefficient.}{%
\fivechoices{+0.50}{+0.25}{0}{-0.25}{-0.50}}{%
\(r_p=(Z_2-Z_1)/(Z_2+Z_1)=(6-2)/(6+2)=+0.50\).}

\bankq{calc-med-097}{Propagation in Different Media}{applied}{B}{impedance,transmission}{Calculate intensity transmission coefficient.}{%
For the boundary in Question 96, find the normal-incidence intensity transmission coefficient \(T_I=4Z_1Z_2/(Z_1+Z_2)^2\).}{%
\fivechoices{0.50}{0.75}{1.00}{1.50}{3.00}}{%
\(T_I=4(2)(6)/(2+6)^2=48/64=0.75\).}

\bankq{calc-med-098}{Propagation in Different Media}{challenge}{C}{refraction,waves}{Calculate refracted wave angle.}{%
A wave crosses from speed \(v_1=\SI{6.0}{m.s^{-1}}\) into speed \(v_2=\SI{3.0}{m.s^{-1}}\) at incidence angle \(30^\circ\). Find the refracted angle.}{%
\fivechoices{\(7.2^\circ\)}{\(12.0^\circ\)}{\(14.5^\circ\)}{\(24.6^\circ\)}{\(30.0^\circ\)}}{%
\(\sin\theta_1/v_1=\sin\theta_2/v_2\), so \(\sin\theta_2=(3/6)\sin30^\circ=0.25\). Thus \(\theta_2=14.5^\circ\).}

\bankq{calc-med-099}{Propagation in Different Media}{challenge}{D}{travel-time,multimedia}{Calculate travel time through layered media.}{%
A pulse travels \(\SI{0.400}{m}\) through water at \(\SI{1500}{m.s^{-1}}\), then \(\SI{1.20}{m}\) through steel at \(\SI{5000}{m.s^{-1}}\). Ignore interface delay. Find the total time.}{%
\fivechoices{\(\SI{0.267}{ms}\)}{\(\SI{0.400}{ms}\)}{\(\SI{0.480}{ms}\)}{\(\SI{0.507}{ms}\)}{\(\SI{0.800}{ms}\)}}{%
\(t=0.400/1500+1.20/5000=2.667\times10^{-4}+2.400\times10^{-4}=5.067\times10^{-4}\,\si{s}=\SI{0.507}{ms}\).}

\bankq{calc-med-100}{Propagation in Different Media}{challenge}{E}{refraction,critical-angle}{Calculate a critical angle for waves.}{%
A wave travels from a medium where its speed is \(\SI{1500}{m.s^{-1}}\) toward one where its speed is \(\SI{3000}{m.s^{-1}}\). Find the critical incidence angle.}{%
\fivechoices{\(15.0^\circ\)}{\(20.0^\circ\)}{\(24.6^\circ\)}{\(26.6^\circ\)}{\(30.0^\circ\)}}{%
At the critical angle, \(\sin\theta_c/v_1=1/v_2\), so \(\sin\theta_c=v_1/v_2=1500/3000=0.5\). Hence \(\theta_c=30.0^\circ\).}

\end{quizbank}

\quizselectall
\let\storedquizmarks\quizmarks
\let\storedquizdifficulty\quizdifficulty
\renewcommand{\quizmarks}[1]{}
\renewcommand{\quizdifficulty}[1]{}
\printquizquestions[2]
\let\quizmarks\storedquizmarks
\let\quizdifficulty\storedquizdifficulty

\begin{quizanswerkeycontent}
  \clearpage
  \begin{answerkey}
    \centering
    \renewcommand{\arraystretch}{1.25}
    \begin{tabular}{*{5}{>{\centering\arraybackslash}p{0.17\linewidth}}}
      \toprule
      \textbf{1.} A & \textbf{2.} B & \textbf{3.} C & \textbf{4.} D & \textbf{5.} E \\
      \textbf{6.} A & \textbf{7.} B & \textbf{8.} C & \textbf{9.} D & \textbf{10.} E \\
      \textbf{11.} A & \textbf{12.} B & \textbf{13.} C & \textbf{14.} D & \textbf{15.} E \\
      \textbf{16.} A & \textbf{17.} B & \textbf{18.} C & \textbf{19.} D & \textbf{20.} E \\
      \textbf{21.} A & \textbf{22.} B & \textbf{23.} C & \textbf{24.} D & \textbf{25.} E \\
      \textbf{26.} A & \textbf{27.} B & \textbf{28.} C & \textbf{29.} D & \textbf{30.} E \\
      \textbf{31.} A & \textbf{32.} B & \textbf{33.} C & \textbf{34.} D & \textbf{35.} E \\
      \textbf{36.} A & \textbf{37.} B & \textbf{38.} C & \textbf{39.} D & \textbf{40.} E \\
      \textbf{41.} A & \textbf{42.} B & \textbf{43.} C & \textbf{44.} D & \textbf{45.} E \\
      \textbf{46.} A & \textbf{47.} B & \textbf{48.} C & \textbf{49.} D & \textbf{50.} E \\
      \textbf{51.} A & \textbf{52.} B & \textbf{53.} C & \textbf{54.} D & \textbf{55.} E \\
      \textbf{56.} A & \textbf{57.} B & \textbf{58.} C & \textbf{59.} D & \textbf{60.} E \\
      \textbf{61.} A & \textbf{62.} B & \textbf{63.} C & \textbf{64.} D & \textbf{65.} E \\
      \textbf{66.} A & \textbf{67.} B & \textbf{68.} C & \textbf{69.} D & \textbf{70.} E \\
      \textbf{71.} A & \textbf{72.} B & \textbf{73.} C & \textbf{74.} D & \textbf{75.} E \\
      \textbf{76.} A & \textbf{77.} B & \textbf{78.} C & \textbf{79.} D & \textbf{80.} E \\
      \textbf{81.} A & \textbf{82.} B & \textbf{83.} C & \textbf{84.} D & \textbf{85.} E \\
      \textbf{86.} A & \textbf{87.} B & \textbf{88.} C & \textbf{89.} D & \textbf{90.} E \\
      \textbf{91.} A & \textbf{92.} B & \textbf{93.} C & \textbf{94.} D & \textbf{95.} E \\
      \textbf{96.} A & \textbf{97.} B & \textbf{98.} C & \textbf{99.} D & \textbf{100.} E \\
      \bottomrule
    \end{tabular}
    \medskip

    \textbf{Total marks: 100.}
  \end{answerkey}
\end{quizanswerkeycontent}

\begin{quizsolutioncontent}
  \clearpage
  \printquizsolutions
\end{quizsolutioncontent}

\end{document}
